% Options for packages loaded elsewhere
\PassOptionsToPackage{unicode}{hyperref}
\PassOptionsToPackage{hyphens}{url}
\documentclass[
  10pt,
  a4paper,
]{article}
\usepackage{xcolor}
\usepackage[margin=22mm]{geometry}
\usepackage{amsmath,amssymb}
\setcounter{secnumdepth}{-\maxdimen} % remove section numbering
\usepackage{iftex}
\ifPDFTeX
  \usepackage[T1]{fontenc}
  \usepackage[utf8]{inputenc}
  \usepackage{textcomp} % provide euro and other symbols
\else % if luatex or xetex
  \usepackage{unicode-math} % this also loads fontspec
  \defaultfontfeatures{Scale=MatchLowercase}
  \defaultfontfeatures[\rmfamily]{Ligatures=TeX,Scale=1}
\fi
\usepackage{lmodern}
\ifPDFTeX\else
  % xetex/luatex font selection
\fi
% Use upquote if available, for straight quotes in verbatim environments
\IfFileExists{upquote.sty}{\usepackage{upquote}}{}
\IfFileExists{microtype.sty}{% use microtype if available
  \usepackage[]{microtype}
  \UseMicrotypeSet[protrusion]{basicmath} % disable protrusion for tt fonts
}{}
\makeatletter
\@ifundefined{KOMAClassName}{% if non-KOMA class
  \IfFileExists{parskip.sty}{%
    \usepackage{parskip}
  }{% else
    \setlength{\parindent}{0pt}
    \setlength{\parskip}{6pt plus 2pt minus 1pt}}
}{% if KOMA class
  \KOMAoptions{parskip=half}}
\makeatother
\setlength{\emergencystretch}{3em} % prevent overfull lines
\providecommand{\tightlist}{%
  \setlength{\itemsep}{0pt}\setlength{\parskip}{0pt}}
\usepackage{mathrsfs}
\usepackage{mathtools}
\usepackage{xurl}
\emergencystretch=3em
\setcounter{tocdepth}{1}
\newcommand{\Beta}{\mathrm{B}}
\DeclareUnicodeCharacter{00B2}{\ensuremath{^2}}
\DeclareUnicodeCharacter{00B1}{\ensuremath{\pm}}
\DeclareUnicodeCharacter{00D7}{\ensuremath{\times}}
\DeclareUnicodeCharacter{2192}{\ensuremath{\to}}
\DeclareUnicodeCharacter{21A6}{\ensuremath{\mapsto}}
\DeclareUnicodeCharacter{2208}{\ensuremath{\in}}
\DeclareUnicodeCharacter{2209}{\ensuremath{\notin}}
\DeclareUnicodeCharacter{2212}{\ensuremath{-}}
\DeclareUnicodeCharacter{2227}{\ensuremath{\wedge}}
\DeclareUnicodeCharacter{2228}{\ensuremath{\vee}}
\DeclareUnicodeCharacter{2260}{\ensuremath{\ne}}
\DeclareUnicodeCharacter{2264}{\ensuremath{\le}}
\DeclareUnicodeCharacter{2265}{\ensuremath{\ge}}
\DeclareUnicodeCharacter{2297}{\ensuremath{\otimes}}
\usepackage{bookmark}
\IfFileExists{xurl.sty}{\usepackage{xurl}}{} % add URL line breaks if available
\urlstyle{same}
\hypersetup{
  hidelinks,
  pdfcreator={LaTeX via pandoc}}

\author{}
\date{}

\begin{document}

{
\setcounter{tocdepth}{3}
\tableofcontents
}
\section{Tau residue norms and the exact change in the Weil
formula}\label{tau-residue-norms-and-the-exact-change-in-the-weil-formula}

The original semiring is \(S=G(\mathbb Z)=\{\tau\}\sqcup\mathbb Z\),
with the original integer operations, \(\tau+x=x\), and \(\tau x=\tau\).
Write \(e=0_{\mathbb Z}\), retaining \(e\ne\tau\). This calculation
starts from its residue and support maps. It does not replace a chosen
residue norm by a different one without recording the map.

The originating programme's \texttt{globalization\_\allowbreak{}note.\allowbreak{}tex},
ideal/congruence/zeta sections, proves the original arithmetic ideal
norm and its zeta function. Its prime spectrum is two-dimensional. The
companion TAU\_PRIME\_SPECTRUM\_DERIVATION.md checks that
classification, its localizations and the coefficient maps in full. Here
the calculation continues from the different residue constructions to
their exact prime distributions and Weil forms.

\subsection{TN1. Three numerical constructions, with their quotient
maps}\label{tn1.-three-numerical-constructions-with-their-quotient-maps}

For \(n\ge1\), reduction of supported integers defines the
support-preserving semiring map \[
\pi_n:S\longrightarrow G(\mathbb Z/n\mathbb Z),\qquad
\tau\longmapsto\tau_n,\quad a\longmapsto(a\bmod n)^\bullet.
\tag{TN1}
\] It is onto and preserves both operations by their definitions. Its
equivalence classes keep \(\tau\) separate from every integer. Its
global-zero fibre is only \(\{\tau\}\). Arithmetic projection gives a
second surjection \[
G(\mathbb Z/n\mathbb Z)\longrightarrow\mathbb Z/n\mathbb Z,
\quad\tau_n\longmapsto0,\quad b^\bullet\longmapsto b.
\tag{TN2}
\] The composite has zero fibre \(I_n=\{\tau\}\cup n\mathbb Z\). The two
target cardinalities are \(n+1\) and \(n\), respectively. At \(n=1\) the
first target is the two-element Boolean semiring and the second is the
one-element ring; this case is retained.

The usual ideal product is \(I_mI_n=I_{mn}\). Every finite sum of
products of multiples of \(m,n\) is a multiple of \(mn\), and every
multiple of \(mn\) is itself such a product; the adjoined identity
\(\tau\) causes no further integer values. Hence the second cardinality
is multiplicative, whereas the first has the exact defect \[
(m+1)(n+1)-(mn+1)=m+n.
\tag{TN3}
\] Consequently the two actual ideal sums on \(I_n\), for \(\Re s>1\),
are \[
Z_{\rm ar}(s)=\sum_{n\ge1}n^{-s}=\zeta(s),\qquad
Z_{\rm card}(s)=\sum_{n\ge1}(n+1)^{-s}=\zeta(s)-1.
\tag{TN4}
\] The second equality is an index change in an absolutely convergent
series. It does not have the prime Euler product obtained by replacing
every \(p\) by \(p+1\).

There is also a precisely defined multiplicative prime weight \[
N_{\rm pr}(n)=\prod_{p\mid n}(p+1)^{v_p(n)},\qquad N_{\rm pr}(1)=1.
\tag{TN5}
\] This agrees with the first residue cardinality at prime integers, but
not at general integers or at \(1\). Unique factorization gives its
separate Euler product \[
E_{\rm pr}(s)=\sum_{n\ge1}N_{\rm pr}(n)^{-s}
=\prod_p\frac1{1-(p+1)^{-s}},\qquad \Re s>1.
\tag{TN6}
\] Absolute convergence follows from \(N_{\rm pr}(n)\ge n\); expanding
finite products and then passing to the absolutely convergent limit
proves the identity. Thus all three functions in TN4 and TN6 are
defined, and none is substituted for another.

\subsection{TN2. The exact mixed-support reason for
TN3}\label{tn2.-the-exact-mixed-support-reason-for-tn3}

For commutative rings \(R,T\), let \(\chi_R:G(R)\to\mathbb B\) send
\(\tau\) to \(0\) and every supported element to \(1\); similarly for
\(T\). There is a semiring isomorphism \[
G(R\times T)\xrightarrow{\sim}
G(R)\times_{\mathbb B}G(T),
\quad \tau\longmapsto(\tau_R,\tau_T),\quad
(r,t)^\bullet\longmapsto(r^\bullet,t^\bullet).
\tag{TN7}
\] The fibre product contains exactly the displayed two types, because
equality of support means either both entries are absent or both are
supported. The map is a bijection. Addition and multiplication on two
supported entries are the ring product operations; entries involving the
absent pair obey the external identity and absorber laws. These checks
prove preservation of both operations, both identities, and the inverse
laws. They also prove the pullback universal property: a pair of
semiring maps with equal support takes values in exactly this subset,
and therefore factors uniquely through it.

If \(m,n\) are coprime, choose integers \(u,v\) with \(um+vn=1\). The
usual reduction map \(\mathbb Z/mn\to\mathbb Z/m\times\mathbb Z/n\) is
injective because divisibility by both coprime integers implies
divisibility by their product; it is surjective because \(bvn+cum\) has
residues \(b,c\). Applying TN7 gives \[
G(\mathbb Z/mn)\cong
G(\mathbb Z/m)\times_{\mathbb B}G(\mathbb Z/n).
\tag{TN8}
\] Its embedding into the full product omits precisely \[
(\mathbb Z/m)^\bullet\times\{\tau_n\},\qquad
\{\tau_m\}\times(\mathbb Z/n)^\bullet.
\tag{TN9}
\] These disjoint mixed-support subsets have sizes \(m,n\). TN3 is
therefore the count of explicitly identified subsets. No supported
cancellation is renamed as absence. In support-separated counts the
vector is \((1,n)\); fibre-product multiplication is componentwise,
\((1,m)(1,n)=(1,mn)\). Forgetting the labels by summing coordinates
gives \(n+1\), and TN3 measures exactly the failure of that sum map to
be multiplicative.

\subsection{TN3. Exact analytic comparison for the prime
weight}\label{tn3.-exact-analytic-comparison-for-the-prime-weight}

On \(\Re s>0\), use the branch \[
\log(1-x^{-s})=-\sum_{j\ge1}\frac{x^{-js}}j,\qquad x\ge2,
\] and define \[
L_\tau(s)=\sum_p\bigl(\log(1-p^{-s})-\log(1-(p+1)^{-s})\bigr),
\quad R_\tau(s)=e^{L_\tau(s)}.
\tag{TN10}
\] This definition has no branch ambiguity because \(|x^{-s}|<1\). To
prove convergence, differentiate the summand with respect to the real
variable \(x\): \[
\log(1-p^{-s})-\log(1-(p+1)^{-s})
=-s\int_p^{p+1}\frac{x^{-s-1}}{1-x^{-s}}\,dx.
\tag{TN11}
\] Fix \(a>0\), and put \(d_a=1-2^{-a}>0\). On \(\Re s\ge a\), the
absolute value of the integrand without \(s\) is at most
\(x^{-a-1}/d_a\). The prime intervals are disjoint subsets of
\([2,\infty)\), so \[
|L_\tau(s)|\le \frac{|s|2^{-a}}{a d_a}.
\tag{TN12}
\] The same domination proves local uniform convergence of the series on
the open half-plane. It therefore defines a holomorphic function, and
its exponential is holomorphic and nowhere zero. On \(\Re s>1\), the
absolutely convergent Euler products give \[
E_{\rm pr}(s)=\zeta(s)R_\tau(s).
\tag{TN13}
\] This supplies a meromorphic continuation of the particular function
TN6 to \(\Re s>0\). Its zero and pole divisor there equals that of
\(\zeta\), with multiplicities, because \(R_\tau\) is a holomorphic
unit. This is a statement about TN6 with TN5, not about TN4's function
\(\zeta-1\), and not about an isomorphism of spectra or trace spaces.

The retained logarithmic-derivative correction is \[
r_\tau(s)=\frac{R_\tau'(s)}{R_\tau(s)}
=\sum_p\left(\frac{\log p}{p^s-1}
-\frac{\log(p+1)}{(p+1)^s-1}\right).
\tag{TN14}
\] One may differentiate locally uniformly in TN10 by holomorphy and
Cauchy's formula on slightly larger compact sets. A direct formula,
including a vertical bound, is \[
r_\tau(s)=\sum_p\int_p^{p+1}x^{-s-1}
\left(\frac{s\log x}{(1-x^{-s})^2}-\frac1{1-x^{-s}}\right)dx,
\tag{TN15}
\] \[
|r_\tau(s)|\le
\frac{2^{-a}}{a d_a}
+\frac{\displaystyle |s|2^{-a}\left(\frac{\log2}{a}+\frac1{a^2}\right)}{d_a^2},
\qquad \Re s\ge a.
\tag{TN16}
\] For TN15, differentiate \((\log x)/(x^s-1)\) in \(x\) and integrate
its negative derivative from \(p\) to \(p+1\). The estimates use
\(\int_2^\infty x^{-a-1}dx=2^{-a}/a\) and
\(\int_2^\infty(\log x)x^{-a-1}dx=2^{-a}(\log2/a+1/a^2)\). They prove
absolute local uniform convergence, and show that \(r_\tau\) has at most
linear vertical growth in every fixed strip inside this half-plane. All
quantities in TN10--TN16 are retained; equality of a zero divisor does
not remove this nonzero correction.

\subsection{TN4. The corresponding prime
distributions}\label{tn4.-the-corresponding-prime-distributions}

Use the programme's actual entire test functions from
WEIL\_PACKET\_ANALYTIC\_DERIVATION.md: \[
A(s)=q_P^\#(s)q_Q(s)v(s)^2,\quad
K(t)=A(\tfrac12+it)=\overline{F_P(t)}F_Q(t),\quad
k(u)=\frac1{2\pi}\int_{\mathbb R}K(t)e^{itu}dt.
\tag{TN17}
\] Its proof WA2--WA3 establishes exponential decay of \(A\) in every
fixed vertical strip and \(|k(u)|\le C_R e^{-R|u|}\) for every fixed
\(R>0\). Define the two explicit prime functionals \[
\begin{aligned}
\mathcal P(k)&=\sum_p\sum_{j\ge1}(\log p)p^{-j/2}
\bigl(k(j\log p)+k(-j\log p)\bigr),\\
\mathcal P_\tau(k)&=\sum_p\sum_{j\ge1}\log(p+1)(p+1)^{-j/2}
\bigl(k(j\log(p+1))+k(-j\log(p+1))\bigr).
\end{aligned}
\tag{TN18}
\] Both sums converge absolutely: for \(R=2\), each is dominated by a
constant times \(\sum_p(\log(p+1))(p+1)^{-5/2}/(1-(p+1)^{-5/2})\), or
the analogous expression with \(p\); each is bounded by a convergent sum
over all integers.

The exact difference is \[
\boxed{\mathcal P(k)-\mathcal P_\tau(k)
=\frac1{2\pi}\int_{\mathbb R}K(t)
\,2\Re r_\tau(\tfrac12+it)\,dt.}
\tag{TN19}
\] To prove it, start with \[
\frac1{2\pi i}\int_{\Re s=3/2}
r_\tau(s)\bigl(A(s)+A(1-s)\bigr)ds.
\] On this line, expand both terms in TN14 as absolutely convergent
geometric series and integrate term by term. Fourier inversion and an
entire contour shift give exactly the difference in TN18, with the
factors \(p^{-j/2}\) and \((p+1)^{-j/2}\). Alternatively shift the
integral itself to \(\Re s=1/2\). No poles are crossed because
\(r_\tau\) is holomorphic for positive real part; its vertical bound
TN16 and the strip decay of \(A\) make the horizontal integrals tend to
zero. Reflect \(t\mapsto-t\) in the second summand. Since TN10 has real
coefficients under conjugation, \(r_\tau(\bar s)=\overline{r_\tau(s)}\).
This gives TN19, including its sign and factor \(2\).

\subsection{TN5. Weil's condition with the full tau
correction}\label{tn5.-weils-condition-with-the-full-tau-correction}

Let \[
w_\infty(t)=\Re\psi(\tfrac14+\tfrac{it}2)-\log\pi,
\qquad
w_{\infty,\tau}(t)=w_\infty(t)-2\Re r_\tau(\tfrac12+it).
\tag{TN20}
\] The complete analytic identity WA24 and TN19 give \[
\boxed{B_h([q_P],[q_Q])
=A(0)+A(1)
+\frac1{2\pi}\int_{\mathbb R}\overline{F_P(t)}F_Q(t)
w_{\infty,\tau}(t)dt
-\mathcal P_\tau(k).}
\tag{TN21}
\] Here \(B_h\) is the literal full-jet packet form with its exact
radical, as constructed in WP28--WP29. To check the sign directly, TN19
says \(\mathcal P=\mathcal P_\tau+D\), where \(D\) is the integral of
\(2\Re r_\tau\). Substitution into the original term \(-\mathcal P\)
subtracts this integral, giving TN20. Thus altering the prime norms
alone while keeping the previous infinite-place term does not produce
the same form; TN20 is its required exact correction.

For \(P=Q\), the right side of TN21 is real. Its nonnegativity is
exactly the original packet's Weil positivity condition, since TN21 is
an equality of forms. This calculation does not prove that condition. It
supplies a full new expression of the condition through the separately
defined prime-residue weight, preserving the support cardinality defect
and the correcting distribution. In particular the degree-three negative
representative WP36 of an assumed off-line packet still gives the exact
value \(-2m\) in TN21.

\subsection{TN6. What this does and does not
identify}\label{tn6.-what-this-does-and-does-not-identify}

There are exact maps from the original tau semiring to two different
residue targets, an exact support fibre product explaining the failure
of scalar count multiplicativity, and a complete analytic comparison
from one specified prime-residue Euler product to its Weil functional.
These statements use the original addition and keep its supported zero.
They do not replace the two-dimensional semiring spectrum by the
one-dimensional arithmetic spectrum, or identify either with Connes's
adele-class orbit space.

The further Connes calculation must retain the supported coordinate-zero
hyperplanes, their stabilizers and transverse additive fields, and the
new absence strata. Those are the objects entering Section VI, equations
(5)--(15), of Alain Connes,
\href{https://arxiv.org/abs/math/9811068v1}{\emph{Trace formula in
noncommutative geometry and the zeros of the Riemann zeta function},
arXiv:math/9811068v1}. The independent endpoint and adelic comparison
continues in TAU\_CONNES\_ENDPOINT\_DERIVATION.md and
TAU\_PRIME\_SPECTRUM\_DERIVATION.md. No equality of their Hilbert-space
traces is assumed from the scalar function identity TN13.

\end{document}
